\newpage
\section{Linear Transformations}
\begin{definition}
Let $U$ and $V$, be vector spaces. A \textbf{linear transformation} from $U$ to $V$ is a function $T: U \to V$ satisfying the following:
\begin{enumerate}
    \item $\forall u_1, u_2 \in U, \quad T(u_1 + u_2) = T(u_1) + T(u_2)$ 
    \item $\forall u \in U, \forall c \in \RR, \quad T(cu) = c \cdot T(u)$
\end{enumerate}
\end{definition}

\begin{remark}
\leavevmode
\begin{enumerate}
    \item If $U = V$, then $T$ is called a \textbf{linear operator} on $V$.
    \item Let $T : U \to V$ be a linear transformation
    \begin{enumerate}
        \item $\forall u_1, u_2 \in U, \quad T(u_1 - u_2) = T(u_1) - T(u_2)$
        \item Denote by $\vec{0}_u$ and $\vec{0}_v$ the respective zero vectors of $U$ and $V$. Then $T(\vec{0}_u) = \vec{0}_v$.
    \end{enumerate}
    Therefore, $T(\vec{0}_u + \vec{0}_u) = \vec{0_v} = T(\vec{0}_u)$
    \item $T: U \to V$ is a linear transformation iff $\forall u_1, u_2 \in U, \forall a, b \in \RR$,
    $$T(au_1 + bu_2) = aT(u_1) + bT(u_2)$$
\end{enumerate}
\end{remark}

\begin{example}
Let $U$ and $V$ be vector spaces. The following are linear transformations:
\begin{enumerate}
    \item $T: U \to V$ given by $T(u) = \vec{0}_v, \forall u \in U$.
    \leavevmode
    
    \textit{Solution.} Let $u_1, u_2 \in U$ and $c \in \RR$. Then,
    $$T(u_1 + u_2) = \vec{0}_v$$
    $$T(u_1) + T(u_2) = \vec{0}_v + \vec{0}_v = \vec{0}_v = T(u_1 + u_2)$$
    $$T(cu_1) = \vec{0}_v$$
    $$c \cdot T(u_1) = c\vec{0}_v = \vec{0}_v = T(cu_1)$$
    Therefore, $T$ is a linear transformation.
    \item $\operatorname{id}_u: U \to U$ given by $\operatorname{id}_u(u) = u$
    \leavevmode

    \textit{Solution.} Let $u_1, u_2 \in U$ and $a, b \in \RR$. Then,
    $$\operatorname{id}_u(au_1 + bu_2) = au_1 + bu_2$$
    $$=a\operatorname{id}_u(u_1) + b\operatorname{id}_u(u_2) $$
    Therefore, $\operatorname{id}_u$ is a linear transformation.
\end{enumerate}
\end{example}

\begin{example}
Let $T: \RR^3 \to \RR^2$ be defined by $T(x, y, z) = (x, y)$.

\textit{Solution.} Let $u_1 = (x_1, y_1, z_1), u_2 = (x_2, y_2, z_2) \in \RR^3$ and $a, b \in \RR$. Then,
$$T(au_1 + bu_2) = T(ax_1 + bx_2, ay_1 + by_2, az_1 + bz_2)$$
$$=(ax_1 + bx_2, ay_1 + by_2)$$
$$= a(x_1, y_1) + b(x_2, y_2)$$
$$=aT(x_1, y_1, z_1) + bT(x_2, y_2, z_2)$$
$$= aT(u_1) + bT(u_2)$$
Therefore, $T$ is a linear transformation.
\end{example}

\newpage
\begin{example}
Let $A \in \matmn$ and $T_A: M_{n \times 1}(\RR) \to M_{m \times 1}(\RR)$ be defined by $T_A(x) = Ax$.

\textit{Solution.} Let $x, y \in M_{n \times 1}(\RR)$ and $a, b \in \RR$. Then,
$$T_A(ax + by) = A(ax + by)$$
$$=aAx + bAy$$
$$=aT_A(x) + bT_A(y)$$
Therefore, $T_A$ is a linear transformation.
\end{example}

\begin{example}
Let $n \in \ZZ^+$ and $D: P_n \to P_{n-1}$, be defined by $D(p(x)) = p\prime(x)$

\textit{Solution.} Let $p(x), q(x) \in P_n$ and $c, d \in \RR$. Then
$$D(cp(x) + dq(x)) = \frac{d}{dx}(cp(x) + dq(x))$$
$$=cp'(x) + dq'(x)$$
$$=cD(p(x)) + dD(q(x))$$
Therefore, $D$ is a linear transformation.
\end{example}

\noindent\rule{\textwidth}{0.4pt}

Now, let us consider some examples which uses counterexample to prove that a function is \textbf{NOT} a linear transformation.

\begin{example}
Let $T: \RR \to \RR$ given by $T(x) = |x|$.

\textit{Solution.} Take $u_1 = 2, u_2 = -2$. Then,
$$T(u_1 + u_2) = T(0) = |0| = 0$$
$$T(u_1) + T(u_2) = |-2| + |2| = 4$$
$$\because T(u_1 + u_2) \neq T(u_1) + T(u_2)$$
Therefore, $T$ is not a linear transformation.
\end{example}

\begin{example}
Let $T: \RR^2 \to P_2$ given by $T(a, b) = x^2 + ax + b$.

\textit{Solution.} Let $(a, b) = (6, 7)$ and $c = 2$. Then,
$$T(c(6, 7)) = T(6c, 7c) = T(12, 14) = x^2 + 12x + 14$$
$$cT(6, 7) = 2(x^2 + 6x + 7) = 2x^2 + 12x + 14$$
$$\because T(cu_1) \neq cT(u_1)$$
Therefore, $T$ is not a linear transformation.
\end{example}

\noindent\rule{\textwidth}{0.4pt}

Now suppose we flip it around, and we would like to know $T(u)$ where $T$ is a linear transformation, $U$ is a vector space, and $u \in U$. How can we derive this? It turns out that we can do so by rewriting a vector $u \in U$ as a \textbf{linear combination}! Recall that $u = c_1u_1 + c_2u_2 + \cdots + c_nu_n$ for some $c_1, c_2, \dots, c_n \in \RR$.
$$T(u) = T(c_1u_1 + c_2u_2 + \cdots + c_nu_n)$$
$$=c_1T(u_1) + c_2T(u_2) + \cdots + c_nT(u_n)$$
Therefore, $T(u_1), T(u_2), \dots, T(u_n)$ are sufficient to determine $T(u)$. Then, $T$ is said to be \textbf{completely determined} by $T(u_1), T(u_2), \dots, T(u_n)$.


\newpage
\begin{example}
Let $T: \RR^2 \to \RR^2$ be a linear transformation such that $T(1, 0) = (2, 1)$ and $T(0, 1) = (-3, 2)$. Find $T(x, y)$ for any $x, y \in \RR$.

\textit{Solution.} Let $(x, y) \in \RR^2$. Then,
$$T(x, y) = T(x(1, 0) + y(0, 1))$$
$$=xT(1, 0) + yT(0, 1)$$
$$=x(2, 1) + y(-3, 2)$$
$$=\boxed{(2x - 3y, x + 2y)}$$
\end{example}

\begin{definition}
Let $T: U \to V$ be a linear transformation. 
\begin{enumerate}
    \item The \textbf{kernel} of $T$ is $\ker T = \{u \in U \mid T(u) = \vec{0}_v\} \subseteq U$.
    \item The \textbf{image/range} of $T$ is $\im T = \{T(u) \mid u \in U\} \subseteq V$.
\end{enumerate}
\end{definition}

\begin{remark}
Let $u \in U$ and $v \in V$.
\begin{enumerate}
    \item $u \in \ker T$ iff $T(u) = \vec{0}_v$
    \item $v \in \im T$ iff $v = T(u)$ for some $u \in U$
\end{enumerate}
\end{remark}

\begin{theorem}
Let $T: U \to V$ be a linear transformation. Then $\ker T \leq U$ and $\im T \leq V$.
\end{theorem}

\begin{definition}
Let $T: U \to V$ be a linear transformation with $U$ and $V$ be finite-dimensional.
\begin{enumerate}
    \item The \textbf{rank} of $T$ is $\rank T = \dim(\im T)$
    \item The \textbf{nullity} of $T$ is $\upsilon(T) = \dim(\ker T)$
\end{enumerate}
\end{definition}

\begin{example}
Let $A \in \matmn$ and $T_A: M_{n \times 1}(\RR) \to M_{m \times 1}(\RR)$ defined by $T_A(x) = Ax$.
\begin{enumerate}[label=(\alph*)]
    \item 
    \begin{align*}
    \ker T_A = \{x \in M_{n \times 1}(\RR) \mid T_A(x) = 0_{m \times 1}\} \\
    =\{x \in M_{n \times 1}(\RR) \mid Ax = 0_{m \times 1}\} \\
    =N(A), \quad \text{the null space of $A$} \\
    \therefore \upsilon(T_A) = \dim(\ker T_A) = \dim N(A) = \upsilon(A)
    \end{align*}
    
    
    \item 
    $$\im T_A = \{T_A(x) \mid x \in M_{n \times 1}(\RR)\}$$
    $$=\{Ax \mid x \in M_{n \times 1}(\RR)\}$$
    $$=\{x_1C_1 + x_2C_2 + \cdots + x_nC_n \mid x_1, x_2, \dots, x_n \in \RR \hspace{0.2cm} \text{and} \hspace{0.2cm} C_i = \text{col $i$ of $A$}\}$$
    $$=\operatorname{span}\{C_1, C_2, \dots, C_n\}$$
    $$= C(A), \quad \text{the column space of $A$}$$
    $$\therefore \rank(T_A) = \dim(\im T) = \dim C(A) = \rank A$$
\end{enumerate}
\end{example}

\newpage
\begin{example}
Let $T: \RR^3 \to \RR^3$ defined by $T(x, y, z) = (x + z, x + y, y - z)$.\footnote{As an exercise prove that this is a linear transformation.} Find the bases for the image and kernel of $T$.

\textit{Solution.}

\begin{enumerate}[label=(\alph*)]
    \item Recall that $(x, y, z) \in \ker T$ iff $T(x, y, z) = (0, 0, 0)$. Then,
    $$(x + z, x + y, y - z) = (0, 0, 0) \quad \longleftrightarrow \quad \begin{cases}
        x + z = 0 \\
        x + y = 0 \\
        y - z = 0
    \end{cases}$$
    Turning this into a coefficient matrix,
    $$\bmat{1 & 0 & 1 \\ 1 & 1 & 0 \\ 0 & 1 & -1} \xrightarrow{\text{RREF}} \bmat{1 & 0 & 1 \\ 0 & 1 & -1 \\ 0 & 0 & 0} \longleftrightarrow \begin{cases}
        x + z = 0 \\
        y - z = 0
    \end{cases}$$
    Since the third column is a free variable, we let $z = t$. Then, $x = -t$ and $y = t$. Therefore $\ker T = \{(-t, t, t) \mid t \in \RR\} = \spn\{(-1, 1, 1)\}$. 
    
    Then $B = \{(-1, 1, 1)\}$ is a basis for $\ker T$, and $\upsilon(T) = 1$.

    \item Let us find a basis for $\im T$,
    $$\im T = \{T(x, y, z) \mid x, y, z \in \RR\}$$
    $$=\{(x + z, x + y, y - z) \mid x, y, z \in \RR\}$$
    $$=\{x(1, 1, 0) + y(0, 1, 1) + z(1, 0, -1) \mid x, y, z \in \RR\}$$
    $$=\spn\{(1, 1, 0), (0, 1, 1), (1, 0, -1)\}$$
    Then, let us check whether this is linearly independent. Form,
    $$x(1, 1, 0) + y(0, 1, 1) + z(1, 0, -1) = (0, 0, 0)$$
    The coefficient matrix is,
    $$\bmat{1 & 0 & 1 \\ 1 & 1 & 0 \\ 0 & 1 & -1} \xrightarrow{\text{RREF}} \bmat{1 & 0 & 1 \\ 0 & 1 & -1 \\ 0 & 0 & 0}$$
    Since there is a free variable in the third column, we must delete that. 
    
    Therefore $B = \{(1, 1, 0), (0, 1, 1)\}$ is a basis for $\im T$ and $\rank T = 2$
\end{enumerate}
\end{example}

\leavevmode

\begin{theorem}[Rank-Nullity Theorem]
Let $U$ and $V$ be vector spaces with $\dim U = n < \infty$. If $T: U \to V$ is a linear transformation, then
$$\dim U = \rank T + \upsilon(T)$$
\end{theorem}

\begin{definition}
Let $T: U \to V$ be a linear transformation.
\begin{enumerate}
    \item $T$ is \textbf{one-to-one} if $\forall u_1, u_2 \in U$. $T(u_1) = T(u_2)$ iff $u_1 = u_2$.
    \item $T$ is \textbf{onto} if $\im T = V$, i.e. $\forall v \in V, \hspace{0.2cm} \exists u \in U$ such that $v = T(u)$.
    \item $T$ is a \textbf{bijection/isomorphism} if $T$ is both one-to-one and onto.
\end{enumerate}
\end{definition}

\newpage
\begin{remark}
Let $U$ and $V$ be vector spaces,
\begin{enumerate}
    \item If there exists an isomorphism $T: U \to V$, then $U$ and $V$ are said to be isomorphic, denoted by $U \cong V$.
    \item If $T: U \to V$ is a linear transformation, then $T$ is one-to-one iff $\ker T = \{\vec{0}_v\}$. Only contains the zero vector.
    \item $T$ is onto iff $\rank T = \dim V$.   
\end{enumerate}
\end{remark}

\lline

For the next few examples, we will \textbf{determine the mapping} of a given function using the remarks given above.

\begin{example}
Determine the mapping of a function $\id_U: U \to U$ defined by $\id_U(u) = u$.
    
\textit{Solution.} 

\begin{enumerate}
    \item $u \in \ker(\id_U)$ iff $\id_U(u) = \vec{0}_u$ and $u = \vec{0}_u$. 
    
    Therefore, $\ker(\id_U) = \vec{0}_u$, hence, $\id_U$ is one-to-one.

    \item $\im(\id_U) = \{\id_U(u) \mid u \in U\} = \{u \mid u \in U\}$

    Therefore $\id_U$ is onto. And since $\id_U$ is both one-to-one and onto, it is an isomorphism.
\end{enumerate}
\end{example}

\begin{example}
Determine the mapping of a function $T: \RR^2 \to \RR^3$ defined by $T(x, y) = (x, y, 0)$.

\textit{Solution.}

\begin{enumerate}[label=(\alph*)]
    \item $(x, y) \in \ker T$ iff $T(x, y) = (0, 0, 0)$
    $$(x, y, 0) = (0, 0, 0) \Longrightarrow x = y = 0$$
    Therefore, $\ker T = \{\vec{0}\}$, hence, $T$ is one-to-one.

    \item $\im T = \{T(x, y) \mid x, y \in \RR\}$
    $$=\{(x, y, 0) \mid x, y \in \RR\}$$
    $$=\{x(1, 0, 0) + y(0, 1, 0) \mid x, y \in \RR\}$$
    $$=\spn\{(1, 0, 0), (0, 1, 0)\}$$
    Since neither of them are scalar multiples of each other it is a basis for $\im T$. And since $\rank T \neq \dim \RR^3$, it is not onto.

    Therefore, $T$ is one-to-one.
\end{enumerate}
\end{example}

\begin{example}
Determine the mapping of a function $T:\RR^3 \to P_2$ defined by $T(a, b, c) = (a + b)x^2 + (b + c)x + c$.

\textit{Solution.}

\begin{enumerate}[label=(\alph*)]
    \item $(a, b, c) \in \ker T$ iff $(a + b)x^2 + (b + c)x + c = 0$
    $$\begin{cases}
        a + b = 0 \\
        b + c = 0 \\
        c = 0
    \end{cases} \quad \longleftrightarrow \quad (a, b, c) = (0, 0, 0)$$
    Therefore $\ker T = \{(0, 0, 0)\}$, hence, $T$ is one-to-one.

    \item $\im T = \{T(a, b, c) \mid a, b, c \in \RR\}$
    $$=\{(a + b)x^2 + (b + c)x + c \mid a, b, c \in \RR\}$$
    $$=\{a(x^2) + b(x^2 + x) + c(x + 1) \mid a, b, c \in \RR\}$$
    $$=\spn\{\underbrace{x^2, x^2 + x, x + 1}_B\}$$
    We can verify that $B$ is a basis for $\im T$. And since $\rank T = 3 = \dim P_2$, hence, $T$ is onto.

    Therefore, $T$ is an isomorphism, and $\RR^3 \cong P_2$.
\end{enumerate}
\end{example}

\begin{theorem}
Let $U, V, W$ be vector spaces and $T: U \to V$ and $S: V \to W$ be linear transformations.
\begin{enumerate}
    \item The function $S \circ T: U \to W$ is a linear transformation. If $S$ and $T$ are isomorphic, then so is $S \circ T$.
    \item $T$ is an isomorphism iff there exists a unique function $T\inv: V \to U$ such that,
    $$(T \circ T\inv)(v) = v, \quad \forall v \in V$$
    $$(T\inv \circ T)(u) = u, \quad \forall u \in U$$
\end{enumerate}
\end{theorem}

\begin{remark}
Let $U, V, W$ be vector spaces
\begin{enumerate}
    \item If $T: U \to V$ is an isomorphism, then the function $T\inv: V \to U$ is also an isomorphism, called the inverse of $T$.
    \item
    \begin{enumerate}
        \item $U \cong U$
        \item If $U \cong V$, then $V \cong U$
        \item If $U \cong V$ and $V \cong W$, then $U \cong W$
    \end{enumerate}
\end{enumerate}
\end{remark}

\begin{theorem}
Let $U$ and $V$ be finite-dimensional vector spaces.
\begin{enumerate}
    \item If $\dim V = n$, then $V \cong \RR^n$
    \item $U \cong V$ iff $\dim U = \dim V$
    \item Let $T: U \to V$ be a linear transformation. If $\dim U = \dim V$, then the following are equivalent:
    \begin{enumerate}
        \item $T$ is one-to-one
        \item $T$ is onto
        \item $T$ is an isomorphism
    \end{enumerate}
\end{enumerate}
\end{theorem}

\begin{example}
\leavevmode
\begin{enumerate}
    \item $M_{n \times 1}(\RR) \cong \RR^n$
    \item $\matmn \cong M_{n \times m}(\RR) \cong \RR^{mn}$
    \item $P_n \cong \RR^{n + 1}$
    \item Show that $T: \RR^2 \to \CC$ given by $T(a, b) = (a + b) + (a - b)i$ is an isomorphism.
\end{enumerate}
\end{example}

\begin{remark}
Let $T: U \to V$ be a linear transformation.
\begin{enumerate}
    \item $T$ is not necessarily and isomorphism even when $U \cong V$. It only guarantees that \ul{such functions exist}.
    \item If $\dim U \neq \dim V$, then $T$ is \textbf{NOT} an isomorphism.
    \begin{enumerate}
        \item If $\dim U > \dim V$, then $T$ is not one-to-one.
        \item If $\dim U < \dim V$, then $T$ is not onto.
    \end{enumerate}
\end{enumerate}
\end{remark}